\chapter{Open questions}\label{ch:open}

\section{Not run}
\begin{itemize}
\item \textbf{CIME.}
The CIME field map is not available to the benchmarks yet.
The node density per harmonic period of the CIME map, and the size of the bilinear excess on it, are therefore unknown.
Case C (acceleration through the dees) does not need it, but waits for a check of the RF time convention of each
code on a one-gap test.
\item \textbf{Extraction (case F).}
The CIME deck has one electrostatic deflector and no transfer-line element, so there is nothing after extraction to
compare between LIONS, Xsuite and Cheetah yet.
\end{itemize}

\section{Untested hypotheses}
\begin{itemize}
\item \textbf{Kink noise on the bilinear path.}
The bilinear one-turn map is only piecewise smooth, and probes that straddle a cell edge mix two slopes.
Test: a Jacobian step scan with the closed orbit moved across a cell edge.
\item \textbf{Residual on RF-Track with the bilinear field.}
RF-Track on the bilinear PSI field still leaves $|\det M_r - 1|$ up to \GTwoCRfMax{} (\cref{tab:G2det}).
Adaptive Runge--Kutta also loses order at gradient jumps, and the section crossing adds its own error.
These were not separated.
\item \textbf{OPAL offsets.}
In B0, OPAL has an energy-dependent closed-orbit offset of order $10^{-8}$ shared by LF2 and RK4.
In A, OPAL RK4 has a phase error floor near \ARkFourPhaseFloor.
Neither is explained.
\item \textbf{Gradients on F1.}
Whether the Hermite field closes the gap between autograd and finite differences on measured maps has not been measured.
\end{itemize}

\section{Physics questions for the LIONS authors}
Two behaviours were found while porting and left unchanged pending a decision:
\begin{itemize}
\item A difference between the codes: the Python port adds a vertical bunch spread for \texttt{espace='6d'}, while the
Fortran does so only for \texttt{'3d'} (\texttt{lions\_2.0\_ep.f90}, line 498).
\item Shared by both codes: the correlation code indexes the mass array by family number.
This is wrong for families after the first.
\end{itemize}

\section{Lessons}
\begin{enumerate}
\item The field smoothness sets the error level more than the integrator does.
On the bilinear maps of case H, neither Boris nor RK4 meets the $10^{-4}$ target at any step count up to 2880 per turn
(\cref{tab:Hdecision}).
A $C^1$ interpolation at \HCostHerm{} times the cost lowers $|\det M-1|$ by \HDetRatioMin{} to \HDetRatioMax{} times.
\item Count nodes per harmonic period, not degrees.
\item Fixed-step integration across gradient jumps breaks area preservation of the one-turn map by orders of magnitude
more than on a smooth field, with Boris and with RK4.
\item Every convention (units, signs, direction, reference azimuth) should be settled by a two-reading run.
\item A criterion built by the analysis that it judges can be tautological.
Write the triangle inequality down before using one.
\end{enumerate}
